\section{Introduction}

The intricacies and fluctuations inherent in financial markets pose significant challenges for making high-quality, sequential investment decisions. In tasks such as single stock trading and portfolio management, each intelligent decision is driven by multiple market interactions and the integration of diverse information streams, characterized by varying levels of timeliness and modalities \cite{markowitz1952portfolio, dang2019squawk}. The primary objective of these tasks is to maximize profit while managing present market risks in an open-ended environment.

In practice, trading firms often depend on synthesized teamwork, structured hierarchically with functional roles such as data analysts, risk analysts, and portfolio managers communicating across levels \cite{maginn2007managing, radner1993organization}. These roles are responsible for the careful integration of diverse resources. However, the cognitive limitations of human team members can hinder their capacity to rapidly process market signals and achieve optimal investment outcomes \cite{miller1956magical}.

%To improve investment returns and overcome human limitations, studies have employed methods such as deep reinforcement learning (DRL) to develop agent systems that simulate market environments and automate investment strategies \cite{wang2021commission, han2023select, qin2024earnhft}. Meanwhile, advancements in large language models (LLMs) demonstrate potential in complex tasks such as reasoning \cite{huang2022towards, jin2024impact}, tool-using \cite{cai2023large}, planning \cite{ruan2023tptu}, decision-making \cite{song2023llm, paranjape2023art}, and even in several financial application \cite{NEURIPS2023_6a386d70, 10.1145/3637528.3671554, xie2024finben, hu2024no, yang2024ucfe}, promising to surpass existing agent architectures. Language agents are known for human-like communication and adaptable prompt-based structures tailored for diverse decision-making settings \cite{park2023generative, wu2023autogen, qiao2024autoact, hong2023metagpt}. For optimal decision-making performance, two key elements need to be considered: 1) Organizing agents to promote effective teamwork and efficient communication, and 2) Enabling agents to continually learn and refine their actions. Studies indicate that imitating human organizational structures can effectively coordinate language agents for specific tasks \cite{guo2024embodied, li2024agent, chen2023agentverse}. Additionally, recent developments in textual gradient-based prompt optimization \cite{pryzant2023automatic, tang2024unleashing} and verbal reinforcement \cite{yao2023retroformer,shinn2024reflexion} have proven to be practical for enhancing the reasoning and decision-making abilities of language agents iteratively.
To enhance investment returns and address the limitations of human decision-making, various studies have explored methods such as deep reinforcement learning (DRL) to develop agent systems that simulate market environments and automate investment strategies \cite{wang2021commission, han2023select, qin2024earnhft}. Concurrently, advancements in large language models (LLMs) have shown great potential in performing complex tasks, including reasoning \cite{huang2022towards, jin2024impact}, tool-using \cite{cai2023large}, planning \cite{ruan2023tptu}, decision-making \cite{song2023llm, paranjape2023art}, and even in various financial applications \cite{NEURIPS2023_6a386d70, 10.1145/3637528.3671554, xie2024finben, hu2024no, yang2024ucfe}, suggesting they may surpass existing agent architectures. Language agents, in particular, are distinguished by their human-like communication and flexible, prompt-based structures, making them well-suited to diverse decision-making settings \cite{park2023generative, wu2023autogen, qiao2024autoact, hong2023metagpt}.

To achieve optimal decision-making performance, two critical factors must be considered: (1) Organizing agents to facilitate effective teamwork and efficient communication, and (2) Enabling agents to continuously learn and refine their actions. Studies have shown that mimicking human organizational structures can successfully coordinate language agents for specific tasks \cite{guo2024embodied, li2024agent, chen2023agentverse}. Additionally, recent advances in textual gradient-based prompt optimization \cite{pryzant2023automatic, tang2024unleashing} and verbal reinforcement \cite{yao2023retroformer, shinn2024reflexion} have proven effective in iteratively improving the reasoning and decision-making capabilities of language agents.

Language agent systems designed for financial decision-making, such as \textsc{FinGPT} \cite{yang2023fingpt}, \textsc{FinMem} \cite{yu2023finmem}, and \textsc{FinAgent} \cite{zhang2024finagent}, have shown strong performance. However, they face several limitations. First, their reliance on agents' risk preferences based on short-term market fluctuations fails to control long-term risk exposure, potentially overlooking fundamental factors driving investment returns. A more effective approach is to quantify investment risks using established \textbf{measures of risk} from quantitative finance \cite{delbaen2000coherent, fabozzi2010quantitative}. Second, these systems are often limited to single-asset trading tasks, making them less adaptable to multi-asset financial applications like portfolio management. Third, they place significant pressure on a single agent to understand and process information within a constrained context window, which can degrade decision quality. Although approaches like \textsc{StockAgent} \cite{liuai} use multi-agent systems for stock trading, their reliance on extensive discussions between numerous LLM agents leads to high communication costs and slow decision-making. Moreover, the absence of a clear optimization objective can compromise outcome effectiveness. Additional related work in the literature is discussed in the Appendix \ref{related_work}.

To address these issues, we propose \textsc{FinCon}, an LLM-based multi-agent framework for critical financial tasks, such as single-stock trading and portfolio management, as shown in Figure~\ref{fig0:framework_overview}. Our main contributions are: \textbf{1)} Inspired by real-world investment roles, we introduce a novel \textbf{Synthesized Manager-Analyst hierarchical communication structure} with a risk-control component. This structure allocates financial data from different sources to corresponding functional analyst agents, allowing them to focus on specific insights, while the manager consolidates these inputs to make informed trading decisions. The streamlined communication reduces redundant peer-to-peer interaction, lowering costs and improving efficiency. \textbf{2)} Our framework generalizes beyond stock trading to handle \textbf{portfolio management}, an area not previously addressed by other financial language agent systems. \textbf{3)} We developed a \textbf{dual-level risk control component} to update risk assessments both within and across episodes. Within episodes, risk is supervised using the Conditional Value at Risk (CVaR), a quantile-based risk measure \cite{kuester2006value}. Across episodes, we introduced a \textbf{verbal reinforcement} mechanism, where investment beliefs are updated based on reasoning trajectories and profit-and-loss (PnL) trends, distilled into \textbf{conceptual perspectives}. These insights are selectively back-propagated from the manager to relevant analyst agents. Our ablation studies demonstrate the effectiveness of this risk control design in managing market risk and enhancing trading performance. 

\begin{figure}[ht]
\centering
\includegraphics[height=3.5cm,width=1.0\textwidth]{figure/FinCon_framework_v4.jpg}
\caption{\footnotesize{The general framework of \textsc{FinCon}.}}
\label{fig0:framework_overview}
\end{figure}

%Language agent systems adapted for financial decision-making, such as \textsc{FinGPT} \cite{yang2023fingpt}, \textsc{FinMem} \cite{yu2023finmem} and \textsc{FinAgent} \cite{zhang2024finagent}, have demonstrated impressive performance. However, these systems exhibit several limitations. Firstly, their reliance on the risk preferences of agents based on short-term market fluctuations fails to control overall risk exposures over long-term trading periods. This oversight can lead to potential losses due to the dismissal of fundamental factors driving investment profits. A more effective strategy is to quantify investment risks using \textbf{measures of risk} from quantitative finance \cite{delbaen2000coherent,fabozzi2010quantitative}. Secondly, most existing systems are designed for single-asset trading tasks and show limited adaptability to other financial applications involving multiple assets, such as portfolio management. Thirdly, these systems depend on the information understanding and extraction ability of a single agent, placing heavy pressure on the language model to comprehend and process information within a limited context window. This likely degrades decision quality. A multi-agent system that allocates tasks based on data source and modality could enhance performance. Even though approaches like \textsc{StockAgent} \cite{liuai} utilize a multi-agent system for individual stock trading tasks. However, their decision-making processes heavily depend on extensive discussions among numerous LLM agents, leading to high communication costs and extended processing times. Besides, the lack of a clear optimization goal in such systems can compromise the effectiveness of the outcomes.

%In response to these unaddressed issues, we propose \textsc{FinCon}, an LLM-based multi-agent framework for critical financial decision-making tasks such as single stock trading and portfolio management. Our main contributions are as follows: \textbf{1)} Inspired by real-world investment roles, we designed an innovative \textbf{Synthesized Manager-Analyst hierarchical communication structure} and a risk control component. This communication structure distributes financial data from various sources and modalities to corresponding functional analyst agents, allowing them to focus on distilling essential investment insights and indicators from a single source of information. The manager then makes trading decisions by consolidating these insights and interacting with the risk-control component. The manager and analyst in a synchronous hierarchy to effectively enhance overall comprehension and reasoning capabilities. In this communication network, messages are conveyed only to the workers who need to be informed, eliminating redundant peer-to-peer communication and saving communication costs. \textbf{2)} It is a generalized framework that can conduct not only stock trading but also \textbf{portfolio management}. The latter has not been claimed to be tackled by other financial language agent frameworks. \textbf{3)} We innovatively crafted the agents' persona using a \textbf{dual-level risk control component}, enabling comprehensive risk updates both within and across episodes. Within episodes, risk supervision is conducted using a quantile-based risk measure from quantitative finance, the Conditional Value at Risk (CVaR) \cite{kuester2006value}. For across-episode risk control, we designed a \textbf{unique verbal reinforcement} method to update the decision maker's investment beliefs with conceptualized investment insights to guide future decisions. These investment beliefs are distilled from underlying reasoning trajectory-level profit and loss (PnL) changes and summarized as \textbf{conceptual perspectives} based on the functionality of all analyst agents. Investment beliefs can then be selectively back-propagated from the manager to the corresponding functional analyst agent. Ablation study results illustrate the significant effectiveness of our risk control design in managing market risk and improving trading performance.